%% P02 \dash{} Numerical error, stability and computing in log-space
%%
%% Stub. The brief below is the contract for this program: what it must argue,
%% and what it must buy the reader. Writing the program means deleting the
%% \programstub{} block and replacing it with the program. If the brief turns
%% out to be wrong once the mathematics has been worked, change it and record
%% the contradiction in CLAUDE.md under *Resolved questions*.

\program{Numerical error, stability and computing in log-space}
\label{prog:P02}

\programstub{%
Argument: an algorithm can be correct in exact arithmetic and useless in
floating point; catastrophic cancellation is the mechanism, and there is a
small catalogue of fixes. Payoff: why log-sum-exp instead of exponentiating.
The reader computes the exact logit at which a naive softmax overflows
\code{fp32} and \code{fp16}, then derives the max-subtraction trick and
shows it changes nothing mathematically. Also: why cross-entropy is computed
from logits and never from probabilities; why the naive one-pass variance
formula gives a negative variance and Welford's does not; why summing a
million small gradients in the wrong order loses them.

\medskip
\textbf{Planned length:} 50 frames.
}
